% newchet_subscripts_test.tex
% Test file for the subscript placement of newchet.sty with the stix2 option.  Each pair of
% lines sets the same formulas, first in the math italic or calligraphic font of the stix2
% package and then in newchet's virtual font.  Only the subscripts should move.
% Superscripts, primes, accents and letters side by side keep their places.  The first
% paragraph names the fonts in use, which are newchet-stix2-mathit and newchet-stix2-mathcal
% when newchet finds its virtual fonts.
% Compile with: make newchet-subscripts-test
% or, from this folder: TEXINPUTS=../newchet: TFMFONTS=../newchet: VFFONTS=../newchet: latexmk -pdf newchet_subscripts_test
% Without the three variables, it tests the installed newchet.sty and fonts.

\documentclass[11pt]{article}
\usepackage[stix2,nobib,notitlepage,notoc]{newchet}

% The math italic of the stix2 package, under a family name of its own
\DeclareFontFamily{LS1}{stix2orig}{\skewchar\font127 }
\DeclareFontShape{LS1}{stix2orig}{m}{it}{<-> stix2-mathit}{}
\DeclareFontShape{LS1}{stix2orig}{b}{it}{<-> stix2-mathit-bold}{}
\DeclareMathAlphabet{\stixit}{LS1}{stix2orig}{m}{it}
\SetMathAlphabet{\stixit}{bold}{LS1}{stix2orig}{b}{it}
% The calligraphic capitals of the stix2 package, which have no bold
\DeclareFontFamily{LS2}{stix2calorig}{\skewchar\font127 }
\DeclareFontShape{LS2}{stix2calorig}{m}{n}{<-> stix2-mathcal}{}
\DeclareMathAlphabet{\stixcal}{LS2}{stix2calorig}{m}{n}

% The formulas, for a letter #1
\newcommand*{\patsub}[1]{#1_A}
\newcommand*{\patsup}[1]{#1^A}
\newcommand*{\patboth}[1]{#1_i^2}
\newcommand*{\patprime}[1]{#1'}
\newcommand*{\pathat}[1]{\hat{#1}}
\newcommand*{\pattwice}[1]{#1#1}
\newcommand*{\patcal}[1]{#1_{\mathcal{O}}}
\makeatletter
% \pair\pat{letters}: the italic letters set with \pat, in the stix2 font and then in
% newchet's.  \calpair does the same for calligraphic capitals.
\newcommand*{\pair}[2]{\pairwith\stixit\@firstofone{#1}{#2}}
\newcommand*{\calpair}[2]{\pairwith\stixcal\mathcal{#1}{#2}}
\newcommand*{\pairwith}[4]{%
  \par\noindent\makebox[4.5em][l]{stix2}\@tfor\next:=#4\do{\(#3{#1{\next}}\)\quad}%
  \par\nopagebreak\noindent\makebox[4.5em][l]{newchet}\@tfor\next:=#4\do{\(#3{#2{\next}}\)\quad}%
  \par\smallskip}
% \shifts\pat{letters}: how much narrower each formula is in newchet's font, in thousandths
% of an em of the font, which is the font size.  \calshifts does the same for calligraphic
% capitals.
\newcommand*{\shifts}[2]{\shiftswith\stixit\@firstofone{#1}{#2}}
\newcommand*{\calshifts}[2]{\shiftswith\stixcal\mathcal{#1}{#2}}
\newcommand*{\shiftswith}[4]{%
  \@tfor\next:=#4\do{%
    \setbox0\hbox{\(#3{#1{\next}}\)}\setbox2\hbox{\(#3{#2{\next}}\)}%
    \next\,\the\numexpr\dimexpr\wd0-\wd2\relax*1000/\dimexpr\f@size\p@\relax\relax\quad}}
\makeatother

\begin{document}
\raggedright

\section*{Subscripts with the stix2 option of newchet}

The math italic in use should be \texttt{newchet-stix2-mathit}, or \texttt{stix2-mathit} with the option \texttt{nosubscriptkerns}.  It is

\noindent\(\hbox{\ttfamily\fontname\textfont1}\) in normal math,

\noindent{\boldmath\(\hbox{\ttfamily\fontname\textfont1}\)} in bold math.

The calligraphic font in use should be \texttt{newchet-stix2-mathcal}, or \texttt{stix2-mathcal} with the option \texttt{nosubscriptkerns}, in normal and bold math.  It is

\noindent\(\hbox{\ttfamily\fontname\textfont\symintegrals}\) in normal math,

\noindent{\boldmath\(\hbox{\ttfamily\fontname\textfont\symintegrals}\)} in bold math.

\subsection*{Subscripts}

\pair\patsub{ABCDEFGHIJKLM}
\pair\patsub{NOPQRSTUVWXYZ}
\pair\patsub{abcdefghijklm}
\pair\patsub{nopqrstuvwxyz}

\subsection*{Bold subscripts}

{\boldmath
\pair\patsub{ABCDEFGHIJKLM}
\pair\patsub{NOPQRSTUVWXYZ}
\pair\patsub{abcdefghijklm}
\pair\patsub{nopqrstuvwxyz}
}

\subsection*{Calligraphic subscripts}

\calpair\patsub{ABCDEFGHIJKLM}
\calpair\patsub{NOPQRSTUVWXYZ}
\calpair\patcal{FINPTVW}
{\boldmath\calpair\patsub{FINPTVW}}

\subsection*{Superscripts, both scripts, primes, accents and pairs}

The superscripts, primes, accents and pairs should not move.  With both scripts, only the subscript should move.

\pair\patsup{FPTVWYfj}
\pair\patboth{FPTVWYfj}
\pair\patprime{FPTVWYfj}
\pair\pathat{FPTVWYfj}
\pair\pattwice{FPTVWYfj}
\calpair\patsup{FINPTVW}
\calpair\patboth{FINPTVW}
\calpair\patprime{FINPTVW}
\calpair\pathat{FINPTVW}
\calpair\pattwice{FINPTVW}

\subsection*{Shifts}

How much each subscript moves to the left, in thousandths of an em.  The letters that do not move show 0.

\noindent\shifts\patsub{ABCDEFGHIJKLMNOPQRSTUVWXYZ}

\noindent\shifts\patsub{abcdefghijklmnopqrstuvwxyz}

\noindent In bold: {\boldmath\shifts\patsub{ABCDEFGHIJKLMNOPQRSTUVWXYZ}}

\noindent In bold: {\boldmath\shifts\patsub{abcdefghijklmnopqrstuvwxyz}}

\noindent Calligraphic: \calshifts\patsub{ABCDEFGHIJKLMNOPQRSTUVWXYZ}

The same for superscripts, primes and pairs, which should all show 0.

\noindent\shifts\patsup{ABCDEFGHIJKLMNOPQRSTUVWXYZabcdefghijklmnopqrstuvwxyz}

\noindent\shifts\patprime{FPTVWYfj}\quad\shifts\pattwice{FPTVWYfj}

\noindent Calligraphic: \calshifts\patsup{ABCDEFGHIJKLMNOPQRSTUVWXYZ}

\noindent\calshifts\patprime{FINPTVW}\quad\calshifts\pattwice{FINPTVW}

\end{document}
